\section{Dimension et profondeur dans les anneaux locaux noethériens}
\label{section:0.16-fr}

\subsection{Dimension d'un anneau}
\label{subsection:0.16.1-fr}

\begin{env}[16.1.1]
\label{0.16.1.1-fr}
On appelle dimension (ou dimension de Krull) d'un anneau $A$ et on note
$\dim A$ la dimension (combinatoire) de son spectre
$\operatorname{Spec}(A)$ (14.2.1)~; comme les parties fermées irréductibles de
$\operatorname{Spec}(A)$ sont les
$V(\mathfrak{p})$, où $\mathfrak{p}$ est un idéal premier de $A$ (Bourbaki,
\emph{Alg.~comm.}, chap.~II, \S~4, no~3, cor.~1 de la prop.~14), il revient au
même de dire que $\dim(A)$ est la borne supérieure des longueurs des chaînes
(14.1.1)
\[
  \mathfrak{p}_0\subset\mathfrak{p}_1\subset\cdots\subset\mathfrak{p}_n
\]
d'idéaux premiers de $A$. Si $A\neq0$, on a $\dim A\geq0$. Un anneau artinien
$\neq0$ (en particulier un corps) est de dimension 0 (Bourbaki,
\emph{Alg.~comm.}, chap.~IV, \S~2, no~5, prop.~7)~; un anneau de Dedekind qui
n'est pas un corps (en particulier $\mathbb{Z}$ ou un anneau $k[T]$ de
polynômes à une indéterminée sur un corps $k$) est de dimension 1 (Bourbaki,
\emph{Alg.~comm.}, chap.~VII, \S~2). Un anneau noethérien n'est pas
nécessairement de dimension finie [30].

Si $(\mathfrak{p}_\alpha)$ est la famille des idéaux premiers minimaux de $A$,
on a (14.1.2.1)
\[
\label{0.16.1.1.1-fr}
  \dim(A)=\sup_\alpha(\dim(A/\mathfrak{p}_\alpha)).
\tag{16.1.1.1}
\]
\end{env}

\begin{env}[16.1.2]
\label{0.16.1.2-fr}
Pour tout idéal $\mathfrak{a}$ de $A$, les idéaux premiers de
$A/\mathfrak{a}$ correspondent biunivoquement aux idéaux premiers de $A$
contenant $\mathfrak{a}$, donc
\[
\label{0.16.1.2.1-fr}
  \dim(A/\mathfrak{a})\leq\dim(A).
\tag{16.1.2.1}
\]
Si en outre $\mathfrak{a}$ n'est contenu dans aucun idéal premier minimal
$\mathfrak{p}_\alpha$ de $A$, toute chaîne d'idéaux premiers de $A$ contenant
$\mathfrak{a}$ peut être complétée par un des $\mathfrak{p}_\alpha$, donc on a
\[
\label{0.16.1.2.2-fr}
  1+\dim(A/\mathfrak{a})\leq\dim(A).
\tag{16.1.2.2}
\]
\end{env}

\begin{env}[16.1.3]
\label{0.16.1.3-fr}
\oldpage[0\textsubscript{IV-1}]{23}
Si $S$ est une partie multiplicative de $A$, les idéaux premiers de $S^{-1}A$
correspondent biunivoquement aux idéaux premiers de $A$ ne rencontrant pas
$S$, donc on a
\[
\label{0.16.1.3.1-fr}
  \dim(S^{-1}A)\leq\dim(A).
\tag{16.1.3.1}
\]
Pour tout idéal premier $\mathfrak{p}$ de $A$, on a, par définition (14.2.1)
\[
\label{0.16.1.3.2-fr}
  \dim(A_{\mathfrak{p}})
  =\operatorname{codim}(V(\mathfrak{p}),\operatorname{Spec}(A)).
\tag{16.1.3.2}
\]
Ce nombre est encore appelé hauteur de l'idéal premier $\mathfrak{p}$ et noté
$\operatorname{ht}(\mathfrak{p})$. Plus généralement, pour tout idéal
$\mathfrak{a}$ de $A$, on appelle hauteur de $\mathfrak{a}$ le nombre
\[
\label{0.16.1.3.3-fr}
  \operatorname{ht}(\mathfrak{a})
  =\operatorname{codim}(V(\mathfrak{a}),\operatorname{Spec}(A)).
\tag{16.1.3.3}
\]
Si $(\mathfrak{p}_\lambda)$ est la famille des idéaux premiers minimaux parmi
ceux qui contiennent $\mathfrak{a}$, on a donc (14.2.1)
\[
\label{0.16.1.3.4-fr}
  \operatorname{ht}(\mathfrak{a})
  =\inf_\lambda(\operatorname{ht}(\mathfrak{p}_\lambda)).
\tag{16.1.3.4}
\]
On a évidemment
\[
\label{0.16.1.3.5-fr}
  \dim(A)=\sup_{\mathfrak{m}}(\dim(A_{\mathfrak{m}}))
\tag{16.1.3.5}
\]
où $\mathfrak{m}$ parcourt l'ensemble des idéaux maximaux de $A$.
\end{env}

\begin{env}[16.1.4]
\label{0.16.1.4-fr}
Pour tout idéal premier $\mathfrak{p}$ de $A$, on a (14.2.2.2)
\[
\label{0.16.1.4.1-fr}
  \dim(A_{\mathfrak{p}})+\dim(A/\mathfrak{p})\leq\dim A.
\tag{16.1.4.1}
\]
On dit que $A$ est caténaire (resp. équidimensionnel, équicodimensionnel,
biéquidimensionnel) lorsque $\operatorname{Spec}(A)$ est caténaire (resp.
équidimensionnel,
équicodimensionnel, biéquidimensionnel). Lorsque $A$ est noethérien et
biéquidimensionnel, les deux membres de (16.1.4.1) sont égaux pour tout idéal
premier de $A$ (14.3.5).

Pour tout idéal $\mathfrak{a}$ de $A$, les idéaux premiers de
$A/\mathfrak{a}$ correspondent biunivoquement aux idéaux premiers de $A$
contenant $\mathfrak{a}$, donc, si $A$ est caténaire, il en est de même de
$A/\mathfrak{a}$. De même, pour toute partie multiplicative $S$ de $A$, les
idéaux premiers de $S^{-1}A$ correspondent biunivoquement aux idéaux premiers
de $A$ ne rencontrant pas $S$, donc si $A$ est caténaire, il en est de même de
$S^{-1}A$.

En outre, comme les idéaux premiers de $A$ contenus dans un idéal premier
$\mathfrak{p}$ correspondent biunivoquement aux idéaux premiers de
$A_{\mathfrak{p}}$, pour que $A$ soit caténaire, il faut et il suffit que
$A_{\mathfrak{p}}$ le soit pour tout $\mathfrak{p}$. Dire que $A$ est
caténaire signifie donc ((14.3.2) et (16.1.3.2)) que pour tout couple d'idéaux
premiers $\mathfrak{p}$, $\mathfrak{q}$ de $A$ tels que
$\mathfrak{q}\subset\mathfrak{p}$, on a
\[
\label{0.16.1.4.2-fr}
  \dim(A_{\mathfrak{q}})
  +\dim(A_{\mathfrak{p}}/\mathfrak{q}A_{\mathfrak{p}})
  =\dim(A_{\mathfrak{p}}).
\tag{16.1.4.2}
\]
Si $A$ est biéquidimensionnel, il en est de même de $A_{\mathfrak{p}}$ pour
tout idéal premier $\mathfrak{p}$ de $A$.
\end{env}

\begin{proposition}[16.1.5]
\label{0.16.1.5-fr}
Soit $\rho:A\to B$ un homomorphisme d'anneaux faisant de $B$ une $A$-algèbre
entière. Alors on a $\dim(B)\leq\dim(A)$~; si de plus $\rho$ est injectif, on a
$\dim(B)=\dim(A)$.
\end{proposition}

\begin{proof}
\oldpage[0\textsubscript{IV-1}]{24}
Si $\mathfrak{n}$ est le noyau de $\rho$, $\rho(A)$ est isomorphe à
$A/\mathfrak{n}$, donc $\dim\rho(A)\leq\dim(A)$ par (16.1.2.1), et comme $B$
est entier sur $\rho(A)$, on peut se borner à démontrer la seconde assertion
lorsque $A\subset B$. Si
$\mathfrak{p}'_0\subset\mathfrak{p}'_1\subset\cdots\subset\mathfrak{p}'_n$
est une chaîne d'idéaux premiers de $B$, les $\mathfrak{p}'_i\cap A$ sont
alors des idéaux premiers distincts dans $A$ (Bourbaki, \emph{Alg.~comm.},
chap.~V, \S~2, no~1, cor.~1 de la prop.~1), donc forment une chaîne d'idéaux
premiers. Inversement, pour toute chaîne
$\mathfrak{p}_0\subset\mathfrak{p}_1\subset\cdots\subset\mathfrak{p}_m$
d'idéaux premiers de $A$, on sait qu'il existe pour chaque $i$ un idéal
premier $\mathfrak{p}'_i$ de $B$ tel que
$\mathfrak{p}'_i\cap A=\mathfrak{p}_i$ et que
$\mathfrak{p}'_i\subset\mathfrak{p}'_{i+1}$ pour $0\leq i\leq m-1$
(\emph{loc.~cit.}, cor.~2 du th.~1). D'où la conclusion.
\end{proof}

\begin{proposition}[16.1.6]
\label{0.16.1.6-fr}
Soient $B$ un anneau intègre, $A$ un sous-anneau local de $B$,
$\mathfrak{m}$ son idéal maximal. Supposons que $A$ soit intégralement clos et
que $B$ soit entier sur $A$~; alors, pour tout idéal maximal $\mathfrak{n}$ de
$B$, on a $\dim(B_{\mathfrak{n}})=\dim(A)$.
\end{proposition}

\begin{proof}
La première partie du raisonnement de (16.1.5) montre que
$\dim(B_{\mathfrak{n}})\leq\dim(A)$. Inversement, considérons une chaîne
d'idéaux premiers
$\mathfrak{p}_0\subset\mathfrak{p}_1\subset\cdots\subset
\mathfrak{p}_k=\mathfrak{m}$ de $A$~; on sait que
$\mathfrak{n}\cap A=\mathfrak{m}$ (Bourbaki, \emph{Alg.~comm.}, chap.~V,
\S~2, no~1, prop.~1) et en vertu des hypothèses faites, on peut appliquer le
second théorème de Cohen-Seidenberg (Bourbaki, \emph{Alg.~comm.}, chap.~V,
\S~2, no~4, th.~3), qui prouve l'existence d'une chaîne d'idéaux premiers
$\mathfrak{p}'_0\subset\mathfrak{p}'_1\subset\cdots\subset
\mathfrak{p}'_k=\mathfrak{n}$ de $B$ telle que
$\mathfrak{p}'_i\cap A=\mathfrak{p}_i$ pour tout $i$~; on en déduit aussitôt
$\dim(B_{\mathfrak{n}})\geq\dim(A)$, d'où la proposition.
\end{proof}

\begin{env}[16.1.7]
\label{0.16.1.7-fr}
Soit $M$ un $A$-module~; on appelle dimension de $M$ et on note
$\dim_A(M)$ ou $\dim(M)$ la dimension de l'anneau $A/\operatorname{Ann}(M)$
(isomorphe à l'anneau des homothéties $A_M$ de $M$). On a donc
$\dim(M^k)=\dim(M)$ pour tout $k>0$. Si $M$ est de type fini, les idéaux
premiers de $A$ contenant $\operatorname{Ann}(M)$ sont exactement ceux qui
appartiennent à $\operatorname{Supp}(M)$ ($0_{\mathrm I}$, 1.7.4) et ce
dernier est fermé dans $\operatorname{Spec}(A)$, donc on a alors
\[
\label{0.16.1.7.1-fr}
  \dim(M)=\dim(\operatorname{Supp}(M))\leq\dim(A).
\tag{16.1.7.1}
\]
Si $M$, $N$ sont deux $A$-modules tels que $N\subset M$, on a
\[
\label{0.16.1.7.2-fr}
  \dim(N)\leq\dim(M),\qquad \dim(M/N)\leq\dim(M).
\tag{16.1.7.2}
\]
Si $S$ est une partie multiplicative de $A$ et si $M$ est de type fini, on a
\[
\label{0.16.1.7.3-fr}
  \dim_{S^{-1}A}(S^{-1}M)\leq\dim_A(M).
\tag{16.1.7.3}
\]
En effet, on a
$\operatorname{Ann}(S^{-1}M)=S^{-1}\operatorname{Ann}(M)$ (Bourbaki,
\emph{Alg.~comm.}, chap.~II, \S~2, no~4) et
$S^{-1}A/S^{-1}\operatorname{Ann}(M)=S^{-1}(A/\operatorname{Ann}(M))$.

Par ailleurs, toute chaîne maximale d'idéaux premiers contenant
$\operatorname{Ann}(M)$ se termine par un idéal maximal $\mathfrak{m}$, donc
sa longueur est égale à celle de la chaîne des idéaux premiers correspondants
de $A_{\mathfrak{m}}$ (qui contiennent
$\operatorname{Ann}(M_{\mathfrak{m}})=(\operatorname{Ann}(M))_{\mathfrak{m}}$)~;
de cette remarque et de (16.1.7.3), on tire
\[
\label{0.16.1.7.4-fr}
  \dim_A(M)=\sup_{\mathfrak{m}}
  (\dim_{A_{\mathfrak{m}}}(M_{\mathfrak{m}}))
\tag{16.1.7.4}
\]
où $\mathfrak{m}$ parcourt l'ensemble des idéaux maximaux de $A$.

On dit que $M$ est caténaire (resp. équidimensionnel, équicodimensionnel,
biéquidimensionnel) lorsque $A/\operatorname{Ann}(M)$ l'est.
\end{env}

\begin{proposition}[16.1.8]
\label{0.16.1.8-fr}
\oldpage[0\textsubscript{IV-1}]{25}
Soient $A$ un anneau, $M$ un $A$-module de type fini, $\mathfrak{p}$ un idéal
premier de $A$. On a alors
\[
\label{0.16.1.8.1-fr}
  \dim_{A_{\mathfrak{p}}}(M_{\mathfrak{p}})
  +\dim_A(M/\mathfrak{p}M)\leq\dim_A(M).
\tag{16.1.8.1}
\]
\end{proposition}

\begin{proof}
On a
$\operatorname{Ann}(M_{\mathfrak{p}})=(\operatorname{Ann}(M))_{\mathfrak{p}}$,
donc les idéaux premiers de $A_{\mathfrak{p}}$ contenant
$\operatorname{Ann}(M_{\mathfrak{p}})$ correspondent biunivoquement aux
idéaux premiers de $A$ contenant $\operatorname{Ann}(M)$ et contenus dans
$\mathfrak{p}$. D'autre part, si $\operatorname{Ann}(M)=\mathfrak{n}$,
$V(\operatorname{Ann}(M/\mathfrak{p}M))=V(\mathfrak{p}+\mathfrak{n})$
($0_{\mathrm I}$, 1.7.5), donc les idéaux premiers de $A$ contenant
$\operatorname{Ann}(M/\mathfrak{p}M)$ contiennent
$\mathfrak{p}+\mathfrak{n}$~; d'où la conclusion.

Cette dernière remarque montre en outre que l'on a
\[
\label{0.16.1.8.2-fr}
  \dim_{A/\mathfrak{p}}(M/\mathfrak{p}M)=\dim_A(M/\mathfrak{p}M).
\tag{16.1.8.2}
\]
\end{proof}

\begin{proposition}[16.1.9]
\label{0.16.1.9-fr}
Soient $A$ un anneau noethérien, $\rho:A\to B$ un homomorphisme d'anneaux,
$M$ un $B$-module tel que $M_{[\rho]}$ soit un $A$-module de type fini. Alors
on a
\[
\label{0.16.1.9.1-fr}
  \dim_A(M_{[\rho]})=\dim_B(M).
\tag{16.1.9.1}
\]
\end{proposition}

\begin{proof}
En effet, l'anneau $B_M$ des homothéties du $B$-module $M$ s'identifie à un
sous-anneau de $C=\operatorname{End}_A(M_{[\rho]})$, et $C$ est un $A$-module
de type fini (Bourbaki, \emph{Alg.~comm.}, chap.~III, \S~3, no~1, lemme~2),
donc il en est de même de $B_M$~; d'ailleurs l'anneau $A_M$ des homothéties du
$A$-module $M_{[\rho]}$ est un sous-anneau de $B_M$, image canonique de $A$
dans $C$, donc $B_M$ est fini sur $A_M$~; la conclusion résulte donc de
(16.1.5).
\end{proof}

\begin{env}[16.1.10]
\label{0.16.1.10-fr}
Soit $A$ un anneau noethérien~; si $M$ est un $A$-module de longueur finie, on
sait (Bourbaki, \emph{Alg.~comm.}, chap.~IV, \S~2, no~5, prop.~7 et cor.~1 de
la prop.~7) que $\operatorname{Supp}(M)$ est un espace fini discret, donc
$\dim(M)=0$ si $M\neq0$~; réciproquement, si $M$ est un $A$-module de type
fini tel que $\dim(M)=0$, tout point de $\operatorname{Supp}(M)$ est fermé,
autrement dit est un idéal maximal de $A$, donc (\emph{loc.~cit.}, prop.~7)
$M$ est de longueur finie.
\end{env}

\subsection[Dimension d'un anneau semi-local noethérien]{Dimension d'un
anneau semi-local noethérien\protect\footnote{Ce numéro reproduit l'exposé
donné par Serre dans un cours au Collège de France en 1958.}}
\label{subsection:0.16.2-fr}

\begin{env}[16.2.1]
\label{0.16.2.1-fr}
Soient $A$ un anneau semi-local noethérien, $\mathfrak{r}$ son radical~;
rappelons qu'un idéal de définition $\mathfrak{q}$ de $A$ est un idéal tel que
$\mathfrak{q}\subset\mathfrak{r}$ et que $\mathfrak{q}$ contienne une
puissance de $\mathfrak{r}$~; il revient au même de dire que les seuls idéaux
premiers contenant $\mathfrak{q}$ sont maximaux, ou que $A/\mathfrak{q}$ est
un anneau artinien. Soit $M$ un $A$-module de type fini~; pour tout entier
$n>0$, $M/\mathfrak{q}^nM$ est un $A$-module de longueur finie, égale à
$\sum_{i=0}^{n-1}\operatorname{long}(\gr_i(M))$, où $\gr_\bullet(M)$ est le
$(A/\mathfrak{q})$-module gradué associé à $M$ muni de la filtration
$\mathfrak{q}$-préadique. On sait qu'il existe un polynôme $Q(n)$ à
coefficients rationnels tel que $\operatorname{long}(\gr_n(M))=Q(n)$ pour $n$
assez grand (III, 2.5.4)~; on en déduit aussitôt qu'il existe un polynôme
$P_{\mathfrak{q}}(M,n)$ en $n$, à coefficients rationnels, tel que
\[
  \operatorname{long}(M/\mathfrak{q}^nM)=P_{\mathfrak{q}}(M,n)
\]
pour $n$ assez grand~; il est clair que ce polynôme est unique et que son
coefficient dominant est $>0$ si $M\neq0$~; on dit que c'est le polynôme de
Hilbert-Samuel de $M$ pour la filtration $\mathfrak{q}$-préadique.
\end{env}

\begin{lemma}[16.2.2.1]
\label{0.16.2.2.1-fr}
\oldpage[0\textsubscript{IV-1}]{26}
Soit $(M_n)$ une filtration $\mathfrak{q}$-bonne de $M$ (Bourbaki,
\emph{Alg.~comm.}, chap.~III, \S~3, no~1)~; alors
$\operatorname{long}(M/M_n)$ est égal, pour $n$ grand, à un polynôme en $n$
dont le degré et le coefficient dominant sont les mêmes que ceux de
$P_{\mathfrak{q}}(M,n)$.
\end{lemma}

\begin{proof}
En effet, il existe $n_0$ tel que $M_{n+1}=\mathfrak{q}M_n$ pour $n\geq n_0$,
d'où les inclusions
$\mathfrak{q}^{n+n_0}M\subset M_{n+n_0}=\mathfrak{q}^nM_{n_0}
\subset\mathfrak{q}^nM\subset M_n$ pour $n$ grand, ce qui entraîne
\[
  P_{\mathfrak{q}}(M,n+n_0)\geq
  \operatorname{long}(M/M_{n+n_0})\geq
  P_{\mathfrak{q}}(M,n)\geq\operatorname{long}(M/M_n)
\]
et par suite la conclusion.
\end{proof}

Si $\mathfrak{q}'$ est un second idéal de définition, on a
$\mathfrak{q}'\supset\mathfrak{q}^m$ pour un entier $m$, donc
$P_{\mathfrak{q}'}(M,n)\leq P_{\mathfrak{q}}(M,mn)$~; par suite le degré de
$P_{\mathfrak{q}}(M,n)$ ne dépend pas de l'idéal de définition
$\mathfrak{q}$ considéré~; on le note $d(M)$.

\begin{lemma}[16.2.2.2]
\label{0.16.2.2.2-fr}
Soit $0\to M'\to M\to M''\to0$ une suite exacte de $A$-modules de type
fini. Alors le polynôme
\[
  P_{\mathfrak{q}}(M,n)-P_{\mathfrak{q}}(M',n)
  -P_{\mathfrak{q}}(M'',n)
\]
a un degré $\leq d(M')-1\leq d(M)-1$.
\end{lemma}

\begin{proof}
En effet, la filtration des
$M'_n=M'\cap\mathfrak{q}^nM'$ est $\mathfrak{q}$-bonne (Bourbaki,
\emph{Alg.~comm.}, chap.~III, \S~3, no~1, prop.~1)~; comme
\[
  \operatorname{long}(M/\mathfrak{q}^nM)
  =\operatorname{long}(M''/\mathfrak{q}^nM'')
  +\operatorname{long}(M'/M'_n),
\]
le lemme résulte aussitôt de (16.2.2.1) appliqué à $M'$.
\end{proof}

\begin{theorem}[16.2.3]
\label{0.16.2.3-fr}
\textup{(Krull-Chevalley-Samuel).}
Soient $A$ un anneau semi-local noethérien, $\mathfrak{r}$ son radical,
$M\neq0$ un $A$-module de type fini, $d(M)$ le degré du polynôme de
Hilbert-Samuel de $M$ (pour un idéal de définition quelconque de $A$)~; soit
d'autre part $s(M)$ la borne inférieure des entiers $n$ tels qu'il existe des
éléments $x_1,\ldots,x_n$ de $\mathfrak{r}$ pour lesquels
$M/(x_1M+\cdots+x_nM)$ soit de longueur finie. Alors on a
\[
  d(M)=s(M)=\dim(M).
\]
\end{theorem}

\begin{lemma}[16.2.3.1]
\label{0.16.2.3.1-fr}
Pour tout $x\in\mathfrak{r}$, soit ${}_xM$ le sous-module de $M$ formé des
éléments annulés par $x$. Alors~:
\begin{enumerate}
  \item[\rm{(i)}] On a $s(M)\leq s(M/xM)+1$.
  \item[\rm{(ii)}] Soient $\mathfrak{p}_i$ les idéaux premiers appartenant à
  $\operatorname{Supp}(M)$ et tels que
  \[
    \dim(A/\mathfrak{p}_i)=\dim(M)\quad(1\leq i\leq m).
  \]
\end{enumerate}
Si $x\notin\mathfrak{p}_i$ pour tout $i$, on a
$\dim(M/xM)\leq\dim(M)-1$.
\begin{enumerate}
  \item[\rm{(iii)}] Si $\mathfrak{q}$ est un idéal de définition de $A$, le
  polynôme
  \[
    P_{\mathfrak{q}}({}_xM,n)-P_{\mathfrak{q}}(M/xM,n)
  \]
  est de degré $\leq d(M)-1$.
\end{enumerate}
\end{lemma}

\begin{proof}
L'assertion (i) est triviale, car si $(x_i)_{1\leq i\leq m}$ est tel que pour
le module $N=M/xM$, $N/(x_1N+\cdots+x_mN)$ soit de longueur finie, il suffit
d'observer que ce module est isomorphe à
$M/(xM+x_1M+\cdots+x_mM)$.

Pour prouver (ii), on remarque que, si $\mathfrak{a}$ est l'annulateur de
$M$, les idéaux premiers qui contiennent l'annulateur de $M/xM$ sont ceux qui
contiennent $\mathfrak{a}+Ax$ ($0_{\mathrm I}$, 1.7.5)~; aucun des
$\mathfrak{p}_i$ ne peut donc contenir $\mathfrak{a}+Ax$ et par définition de
$\dim(M)$ (16.1.7) toute chaîne d'idéaux premiers contenant
$\mathfrak{a}+Ax$ a donc une longueur $\leq\dim(M)-1$.

\oldpage[0\textsubscript{IV-1}]{27}

Enfin, pour démontrer (iii), on note que l'on a deux suites exactes
\[
  0\to{}_xM\to M\xrightarrow{x}xM\to0,\qquad
  0\to xM\to M\to M/xM\to0,
\]
et l'assertion découle aussitôt du lemme (16.2.2.2).
\end{proof}

La démonstration de (16.2.3) se fait en trois étapes~:

\begin{env}[16.2.3.2]
\label{0.16.2.3.2-fr}
On a $\dim(M)\leq d(M)$.
\end{env}

Raisonnons par récurrence sur $d(M)$, le cas $d(M)=0$ étant trivial.
Supposons donc $d(M)\geq1$, et soit $\mathfrak{p}_0\in\operatorname{Supp}(M)$
tel que $\dim(A/\mathfrak{p}_0)=\dim(M)$. Cela entraîne que
$\mathfrak{p}_0$ est un élément \emph{minimal} de $\operatorname{Supp}(M)$,
donc il est associé à $M$ (Bourbaki, \emph{Alg.~comm.}, chap.~IV, \S~1,
no~4, th.~2) et $M$ contient par suite un sous-module $N$ isomorphe à
$A/\mathfrak{p}_0$ (\emph{loc.~cit.}, \S~1, no~1)~; comme $d(M)\geq d(N)$
par (16.2.2.2), on est ramené à prouver que $\dim(N)\leq d(N)$. Soit donc
$\mathfrak{p}_0\subset\mathfrak{p}_1\subset\cdots\subset\mathfrak{p}_n$ une
chaîne d'idéaux premiers distincts dans $A$, et montrons que $n\leq d(N)$.
C'est évident si $n=0$. Sinon, il existe
$x\in\mathfrak{p}_1\cap\mathfrak{r}$ tel que $x\notin\mathfrak{p}_0$, car la
relation $\mathfrak{p}_0\supset\mathfrak{p}_1\cap\mathfrak{r}$ entraînerait
$\mathfrak{p}_0\supset\mathfrak{r}$ puisque $\mathfrak{p}_0$ ne contient pas
$\mathfrak{p}_1$, donc $\mathfrak{p}_0$ serait maximal, ce qui est absurde.
On a $\mathfrak{p}_i\in\operatorname{Supp}(N/xN)$, pour $i\geq1$, donc
$n-1\leq\dim(N/xN)$ (16.1.7). Comme ${}_xN=0$ en vertu du choix de $x$, on a
$d(N/xN)\leq d(N)-1$ par (16.2.3.1, (iii)). L'hypothèse de récurrence implique
alors $d(N/xN)=\dim(N/xN)\geq n-1$, donc $n-1\leq d(N)-1$.

\begin{env}[16.2.3.3]
\label{0.16.2.3.3-fr}
On a $d(M)\leq s(M)$.
\end{env}

Soit en effet $(x_i)_{1\leq i\leq n}$ une famille d'éléments de
$\mathfrak{r}$ telle que, si l'on pose
$\mathfrak{a}=x_1A+\cdots+x_nA$, $M/\mathfrak{a}M$ soit de longueur finie.
Les idéaux associés à $M/\mathfrak{a}M$ sont alors maximaux (Bourbaki,
\emph{Alg.~comm.}, chap.~IV, \S~2, no~5, prop.~7) et comme ce sont les seuls
idéaux premiers contenant
$\mathfrak{q}=\mathfrak{a}+(\mathfrak{r}\cap\operatorname{Ann}(M))$,
$A/\mathfrak{q}$ est artinien (\emph{loc.~cit.}, prop.~9) et $\mathfrak{q}$
est un idéal de définition de $A$. Mais on a
$\mathfrak{q}^mM/\mathfrak{q}^{m+1}M=
\mathfrak{a}^mM/\mathfrak{a}^{m+1}M$, et si $M$ est engendré par $r$ éléments
$z_j$, $\mathfrak{a}^mM/\mathfrak{a}^{m+1}M$ est un $(A/\mathfrak{q})$-module
engendré par les images canoniques des
$x_1^{\nu_1}x_2^{\nu_2}\cdots x_n^{\nu_n}z_j$ où
$\sum_i\nu_i=m$, donc par $r\binom{m+n-1}{n-1}$ éléments~; sa longueur est
par suite $\leq ar\binom{m+n-1}{n-1}$, où $a$ est une constante, et on en
déduit aussitôt que le degré de $P_{\mathfrak{q}}(M,n)$ est $\leq n$~; d'où
$d(M)\leq s(M)$.

\begin{env}[16.2.3.4]
\label{0.16.2.3.4-fr}
On a $s(M)\leq\dim(M)$.
\end{env}

Notons d'abord que $n=\dim(M)$ est fini par (16.2.3.2). Raisonnons par
récurrence sur $n$~; la proposition est évidente pour $n=0$, puisque $M$ est
alors de longueur finie (16.1.10). Supposons $n\geq1$, et soient
$\mathfrak{p}_i$ $(1\leq i\leq m)$ les idéaux premiers de
$\operatorname{Supp}(M)$ tels que $\dim(A/\mathfrak{p}_i)=n$~; la définition
des $\mathfrak{p}_i$ montre que ce sont des éléments minimaux de
$\operatorname{Supp}(M)$, donc en nombre fini (Bourbaki, \emph{Alg.~comm.},
chap.~IV, \S~1, no~4, th.~2). Comme $n\geq1$, les $\mathfrak{p}_i$ ne sont pas
maximaux, et il existe donc $x\in\mathfrak{r}$ tel que
$x\notin\mathfrak{p}_i$ pour tout $i$ (Bourbaki, \emph{Alg.~comm.}, chap.~II,
\S~1, no~1, prop.~2). En vertu de (16.2.3.1), on a
$s(M)\leq s(M/xM)+1$ et $\dim(M)\geq\dim(M/xM)+1$~; l'hypothèse de récurrence
entraîne que $s(M/xM)\leq\dim(M/xM)$, d'où la conclusion.

\begin{corollary}[16.2.4]
\label{0.16.2.4-fr}
Sous les hypothèses de (16.2.3), on a
\begin{equation}
\label{0.16.2.4.1-fr}
\tag{16.2.4.1}
  \dim_{\widehat{A}}(\widehat{M})=\dim_A(M).
\end{equation}
\end{corollary}

\oldpage[0\textsubscript{IV-1}]{28}

\begin{proof}
En effet, $M/\mathfrak{r}^nM$ et
$\widehat{M}/\mathfrak{r}^n\widehat{M}$ sont isomorphes, donc
$d(\widehat{M})=d(M)$.
\end{proof}

\begin{corollary}[16.2.5]
\label{0.16.2.5-fr}
La dimension d'un anneau semi-local noethérien $A$ est finie et égale au
nombre minimum d'éléments du radical de $A$ engendrant un idéal de définition.
\end{corollary}

\begin{proof}
C'est l'égalité $s(M)=\dim(M)$ pour $M=A$.
\end{proof}

En particulier~:

\begin{corollary}[16.2.6]
\label{0.16.2.6-fr}
Soient $A$ un anneau local noethérien, $\mathfrak{m}$ son idéal maximal,
$k=A/\mathfrak{m}$ son corps résiduel~; on a
\begin{equation}
\label{0.16.2.6.1-fr}
\tag{16.2.6.1}
  \dim(A)\leq\operatorname{rg}_k(\mathfrak{m}/\mathfrak{m}^2).
\end{equation}
\end{corollary}

\begin{proof}
En effet, on sait que si les $x_i$ $(1\leq i\leq n)$ sont des éléments de
$\mathfrak{m}$ dont les classes mod.~$\mathfrak{m}^2$ forment une base du
$k$-espace vectoriel $\mathfrak{m}/\mathfrak{m}^2$, les $x_i$ engendrent
$\mathfrak{m}$ (Bourbaki, \emph{Alg.~comm.}, chap.~II, \S~3, no~2, cor.~2 de
la prop.~4).
\end{proof}

\begin{proposition}[16.2.7]
\label{0.16.2.7-fr}
Soient $k$ un corps, $B$ une $k$-algèbre graduée de type fini à degrés
positifs engendrée par ses éléments homogènes de degré 1 et telle que $B_0=k$~;
soit $E$ un $B$-module gradué de type fini, de sorte (III, 2.5.4) que, pour $n$
assez grand, $\operatorname{rg}_k(E_n)=r(n)$ est de la forme
$\Phi(n)-\Phi(n-1)$, où $\Phi$ est un polynôme en $n$ de degré $d$. Alors il
existe dans $B$ une famille $(t_i)_{1\leq i\leq d}$ de $d$ éléments homogènes
de degrés $\geq1$ tels que $E/(\sum_i t_iE)$ soit de rang fini sur $k$.
\end{proposition}

\begin{proof}
Soit $\mathfrak{q}$ l'idéal maximal $\bigoplus_{n\geq1}B_n$ de $B$, et, pour
tout $n$, soit $E'_n$ le sous-$B$-module $\bigoplus_{h\geq n+1}E_h$ de $E$.
Tout élément de $B-\mathfrak{q}$ détermine une homothétie injective dans le
$B$-module artinien $E/E'_n$, et cette homothétie est donc bijective
(Bourbaki, \emph{Alg.}, chap.~VIII, \S~2, no~2, lemme~3), ce qui entraîne que
$(E/E'_n)_{\mathfrak{q}}$ s'identifie à $E/E'_n$~; comme le corps résiduel de
$B_{\mathfrak{q}}$ est $k$, on déduit de l'hypothèse que
$\operatorname{long}_{B_{\mathfrak{q}}}((E/E'_n)_{\mathfrak{q}})$ est un
polynôme en $n$ de degré $d$. D'autre part, comme $E$ est un $B$-module de
type fini et que $B$ est engendré par ses éléments homogènes de degré 1, la
filtration $(E'_n)$ sur $E$ est $\mathfrak{q}$-bonne (Bourbaki,
\emph{Alg.~comm.}, chap.~III, \S~1, no~3, lemme~1), donc la filtration
$((E'_n)_{\mathfrak{q}})$ sur $E_{\mathfrak{q}}$ est
$\mathfrak{q}B_{\mathfrak{q}}$-bonne. Puisque $B_{\mathfrak{q}}$ est un anneau
local noethérien, on conclut de (16.2.3) et de (16.2.2.1) que l'on a
$\dim(E_{\mathfrak{q}})=d$. Raisonnons alors par récurrence sur $d$, le lemme
étant évident si $d=0$. Supposons $d\geq1$, et observons que les idéaux premiers
minimaux $\mathfrak{p}_i$ $(1\leq i\leq r)$ dans
$\operatorname{Ass}(E)$ sont \emph{gradués} (Bourbaki, \emph{Alg.~comm.},
chap.~IV, \S~3, no~1, prop.~1), et que $\mathfrak{q}$ est distinct de la
réunion des $\mathfrak{p}_i$ puisque $d\geq1$~; il y a donc un élément
\emph{homogène} $t_1\in\mathfrak{q}$ n'appartenant à aucun des
$\mathfrak{p}_i$ (Bourbaki, \emph{Alg.~comm.}, chap.~III, \S~1, no~4,
prop.~8). Comme les idéaux premiers $(\mathfrak{p}_i)_{\mathfrak{q}}$ sont les
éléments minimaux de $\operatorname{Ass}(E_{\mathfrak{q}})$ (Bourbaki,
\emph{Alg.~comm.}, chap.~IV, \S~1, no~2, prop.~5), on a
$\dim(E_{\mathfrak{q}}/t_1E_{\mathfrak{q}})=d-1$ (16.3.4). Il suffit alors
d'appliquer l'hypothèse de récurrence au $B$-module gradué $E/t_1E$ pour
obtenir la conclusion.
\end{proof}

\subsection{Systèmes de paramètres dans un anneau local noethérien}
\label{subsection:0.16.3-fr}

\begin{proposition}[16.3.1]
\label{0.16.3.1-fr}
Soient $A$ un anneau noethérien, $\mathfrak{p}$ un idéal premier de $A$, $n$ un
entier. Les conditions suivantes sont équivalentes~:
\begin{enumerate}
  \item[$a)$] $\operatorname{ht}(\mathfrak{p})=
  \dim(A_{\mathfrak{p}})\leq n$~;
  \item[$b)$] Il existe $n$ éléments $x_i$ de $A$ $(1\leq i\leq n)$ tels que
  $\mathfrak{p}$ soit minimal dans l'ensemble des idéaux premiers contenant
  l'idéal $\mathfrak{a}$ engendré par les $x_i$ (autrement dit
  $V(\mathfrak{p})$ est une composante irréductible de $V(\mathfrak{a})$).
\end{enumerate}
\end{proposition}

\begin{proof}
La condition $b)$ entraîne que $\mathfrak{a}A_{\mathfrak{p}}$ est un idéal de
définition de $A_{\mathfrak{p}}$, d'où $a)$ en vertu de (16.2.5).
Inversement, si $a)$ est vérifiée, il existe un idéal de définition
$\mathfrak{b}$ de $A_{\mathfrak{p}}$, engendré par $n$ éléments $x_i/s$, avec
$s\in A-\mathfrak{p}$. En vertu de la correspondance biunivoque entre idéaux
premiers de $A_{\mathfrak{p}}$ et idéaux premiers de $A$ contenus dans
$\mathfrak{p}$, l'idéal $\mathfrak{a}$ de $A$ engendré par les $x_i$ vérifie
$b)$.
\end{proof}

\oldpage[0\textsubscript{IV-1}]{29}

Pour $n=1$, on obtient le cas particulier~:

\begin{corollary}[16.3.2]
\label{0.16.3.2-fr}
\textup{(``Hauptidealsatz'').}
Pour qu'un idéal premier dans un anneau noethérien soit de hauteur $\leq1$, il
faut et il suffit qu'il soit minimal dans l'ensemble des idéaux premiers
contenant un idéal principal convenable.
\end{corollary}

\begin{corollary}[16.3.3]
\label{0.16.3.3-fr}
\textup{(Artin-Tate).}
Soit $A$ un anneau intègre noethérien. Les conditions suivantes sont
équivalentes~:
\begin{enumerate}
  \item[$a)$] $A$ est un anneau semi-local de dimension $\leq1$.
  \item[$b)$] Il existe $f\neq0$ dans $A$ tel que $A_f$ soit un corps.
\end{enumerate}
\end{corollary}

\begin{proof}
La condition $a)$ équivaut à dire que $A$ n'a qu'un nombre fini d'idéaux
premiers $\mathfrak{m}_i\neq0$, et que ces idéaux sont tous maximaux ($0$ étant
un idéal premier). Comme le produit des $\mathfrak{m}_i$ n'est pas nul, puisque
$A$ est intègre, un élément $f\neq0$ de ce produit appartient à tous les
$\mathfrak{m}_i$, donc $(0)$ est le seul idéal premier de l'anneau intègre
$A_f$, autrement dit $A_f$ est un corps. (On notera que cette partie de la
démonstration n'utilise pas le fait que $A$ est noethérien.) Prouvons
maintenant que $b)$ entraîne $a)$~; l'hypothèse $b)$ signifie que tout idéal
premier $\mathfrak{p}\neq0$ de $A$ contient $f$. Soient $\mathfrak{p}_i$
$(1\leq i\leq n)$ les idéaux premiers minimaux parmi ceux qui contiennent $f$
(qui sont en nombre fini puisque $A/fA$ est noethérien)~; il suffit de prouver
que les $\mathfrak{p}_i$ sont des idéaux \emph{maximaux}, car alors tout idéal
premier $\neq0$ contient nécessairement un des $\mathfrak{p}_i$, donc lui est
égal. Supposons donc qu'il existe un idéal maximal $\mathfrak{m}$ contenant un
des $\mathfrak{p}_i$ et distinct des $\mathfrak{p}_i$~; $\mathfrak{m}$ n'est
donc pas contenu dans la réunion des $\mathfrak{p}_i$ (Bourbaki,
\emph{Alg.~comm.}, chap.~II, \S~1, no~1, prop.~2), autrement dit il existe
$g\in\mathfrak{m}$ n'appartenant à aucun des $\mathfrak{p}_i$. Soit
$\mathfrak{q}$ un des idéaux premiers minimaux parmi ceux qui contiennent
$g$~; il résulte du ``Hauptidealsatz'' (16.3.2) que $\mathfrak{q}$ est de
hauteur 1, donc $\neq0$~; il contient par suite $f$, donc aussi un des
$\mathfrak{p}_i$, et comme $g\in\mathfrak{q}$ et que $g$ n'appartient à aucun
des $\mathfrak{p}_i$, $\mathfrak{q}$ est distinct des $\mathfrak{p}_i$~; il
serait donc de hauteur $\geq2$, ce qui est contradictoire.
\end{proof}

\begin{proposition}[16.3.4]
\label{0.16.3.4-fr}
Soient $A$ un anneau semi-local noethérien, $\mathfrak{r}$ son radical, $M$ un
$A$-module de type fini, $\mathfrak{p}_i$ les éléments de
$\operatorname{Supp}(M)$ tels que
$\dim(A/\mathfrak{p}_i)=\dim(M)$. Alors les $\mathfrak{p}_i$ sont des éléments
minimaux de $\operatorname{Ass}(M)$, et l'on a
$\dim(M/\mathfrak{p}_iM)=\dim(M)$. Pour tout $x\in\mathfrak{r}$, on a
\begin{equation}
\label{0.16.3.4.1-fr}
\tag{16.3.4.1}
  \dim(M/xM)\geq\dim(M)-1.
\end{equation}
En outre, pour que les deux membres de (16.3.4.1) soient égaux, il faut et il
suffit que $x$ n'appartienne à aucun des $\mathfrak{p}_i$.
\end{proposition}

\begin{proof}
La première assertion est immédiate, car les $\mathfrak{p}_i$ sont par
définition des éléments minimaux de $\operatorname{Supp}(M)$ (16.1.6), et ces
derniers sont aussi éléments minimaux de $\operatorname{Ass}(M)$ (Bourbaki,
\emph{Alg.~comm.}, chap.~IV, \S~1, no~4, th.~2). En outre,
$\operatorname{Supp}(M/\mathfrak{p}_iM)$ est l'ensemble des idéaux premiers
contenant $\mathfrak{p}_i$ ($0_{\mathrm I}$, 1.7.5), donc
$\dim(M/\mathfrak{p}_iM)=\dim(A/\mathfrak{p}_i)=\dim(M)$.

Posons $N=M/xM$. Si $x_1,\ldots,x_n$ sont des éléments de $\mathfrak{r}$ tels
que $N/(x_1N+\cdots+x_nN)$ soit de longueur finie, cela signifie que
$M/(xM+x_1M+\cdots+x_nM)$ est de longueur finie, d'où l'inégalité (16.3.4.1)
en vertu de (16.2.3).

Le fait que les deux membres de (16.3.4.1) soient égaux lorsque $x$
n'appartient à aucun des $\mathfrak{p}_i$ résulte de (16.3.4.1) et de
(16.2.3.1). Inversement, si
\oldpage[0\textsubscript{IV-1}]{30}
$\dim(M/xM)=\dim(M)-1$, aucun des $\mathfrak{p}_i$ ne peut appartenir à
$\operatorname{Supp}(M/xM)$, et les idéaux premiers de ce support sont ceux
qui contiennent $\operatorname{Ann}(M)+Ax$~; comme les $\mathfrak{p}_i$
contiennent $\operatorname{Ann}(M)$, ils ne peuvent contenir $x$.
\end{proof}

\begin{corollary}[16.3.5]
\label{0.16.3.5-fr}
Avec les notations de (16.3.4), pour tout idéal
$\mathfrak{a}\subset\mathfrak{p}_i$, on a
$\dim(M/\mathfrak{a}M)=\dim(M)$ et si l'on pose
$N=M/\mathfrak{a}M$, la relation $\dim(M/xM)=\dim(M)-1$ entraîne
$\dim(N/xN)=\dim(N)-1$.
\end{corollary}

\begin{proof}
En effet, $\operatorname{Supp}(N)$ est l'ensemble des idéaux premiers
contenant $\mathfrak{a}+\operatorname{Ann}(M)$, donc $\mathfrak{p}_i$
appartient à $\operatorname{Supp}(N)$ et comme
$\dim(N)\leq\dim(M)=\dim(A/\mathfrak{p}_i)$, on a
$\dim(A/\mathfrak{p}_i)=\dim(N)$~; en outre les idéaux premiers
$\mathfrak{q}\in\operatorname{Supp}(N)$ tels que
$\dim(A/\mathfrak{q})=\dim(N)$ sont certains des $\mathfrak{p}_j$~; si $x$
n'appartient à aucun des $\mathfrak{p}_j$, on a donc
$\dim(N/xN)=\dim(N)-1$ en vertu de (16.3.4).
\end{proof}

\begin{definition}[16.3.6]
\label{0.16.3.6-fr}
Soient $A$ un anneau semi-local noethérien, $\mathfrak{r}$ son radical, $M$ un
$A$-module de type fini, et posons $n=\dim(M)$. On appelle \emph{système de
paramètres pour $M$} tout système $(x_i)_{1\leq i\leq n}$ de $n$ éléments de
$\mathfrak{r}$ tel que $M/(x_1M+\cdots+x_nM)$ soit de longueur finie.
\end{definition}

On a vu (16.2.3) qu'il existe toujours de tels systèmes.

On a vu au cours de la démonstration de (16.2.3.3) que si
$\mathfrak{J}=\sum_{i=1}^n x_iA$ est tel que $M/\mathfrak{J}M$ soit de
longueur finie, $\mathfrak{q}=\mathfrak{J}+\operatorname{Ann}(M)$ est un
\emph{idéal de définition} de $A$, et réciproquement, s'il en est ainsi,
$M/\mathfrak{J}M$ est un module de type fini sur l'anneau artinien
$A/\mathfrak{q}$, donc est de longueur finie. Il revient donc au même de dire
que $(x_i)_{1\leq i\leq n}$ est un système de paramètres \emph{pour $M$} ou
\emph{pour $A/\operatorname{Ann}(M)$} (ou encore que leurs images dans
$A/\operatorname{Ann}(M)$ forment un système de paramètres pour cet anneau).
Si $(x_i)_{1\leq i\leq n}$ est un système de paramètres pour $M$, il en est de
même de $(x_i^k)$ pour tout entier $k\geq1$, car on peut se borner, en vertu de
ce qui précède, au cas où $M=A$, et si $\mathfrak{J}$ est un idéal de
définition de $A$, l'idéal $\sum_{i=1}^n x_i^kA$ contient
$\mathfrak{J}^{kn}$, donc est aussi un idéal de définition.

\begin{proposition}[16.3.7]
\label{0.16.3.7-fr}
Soient $A$ un anneau semi-local noethérien, $\mathfrak{r}$ son radical,
$x_1,\ldots,x_k$ des éléments de $\mathfrak{r}$, $M$ un $A$-module de type
fini. On a alors
\begin{equation}
\label{0.16.3.7.1-fr}
\tag{16.3.7.1}
  \dim(M/(x_1M+\cdots+x_kM))\geq\dim(M)-k.
\end{equation}
En outre, pour que les deux membres de (16.3.7.1) soient égaux, il faut et il
suffit qu'il existe des éléments $x_i$
$(k+1\leq i\leq n=\dim(M))$ de $\mathfrak{r}$ tels que
$(x_i)_{1\leq i\leq n}$ soit un système de paramètres pour $M$.
\end{proposition}

\begin{proof}
L'inégalité (16.3.7.1) résulte de (16.3.4.1) par récurrence sur $k$. Si les
deux membres de (16.3.7.1) sont égaux, et si $x_{k+1},\ldots,x_n$ est un
système de paramètres pour $M/(x_1M+\cdots+x_kM)$, il est clair que
$M/(x_1M+\cdots+x_kM+x_{k+1}M+\cdots+x_nM)$ est de longueur finie, donc
$(x_i)_{1\leq i\leq n}$ est un système de paramètres pour $M$~;
réciproquement, s'il existe des $x_i$ $(k+1\leq i\leq n)$ ayant cette
propriété et si l'on pose $N=M/(x_1M+\cdots+x_kM)$, il est clair que
$N/(x_{k+1}N+\cdots+x_nN)$ est de longueur finie, donc
$\dim(N)\leq n-k$, et les deux membres sont égaux en vertu de (16.3.7.1).
\end{proof}

\oldpage[0\textsubscript{IV-1}]{31}

\begin{proposition}[16.3.8]
\label{0.16.3.8-fr}
Soient $A$ un anneau semi-local noethérien, $X=\operatorname{Spec}(A)$ son
spectre, $\mathfrak{a}$ un idéal de $A$ distinct de $A$ tel que
$\operatorname{codim}(V(\mathfrak{a}),X)=r>0$. Il existe alors $r$ éléments
$x_1,\ldots,x_r$ de $\mathfrak{a}$, faisant partie d'un système de paramètres
de $A$, tels que
\[
  \operatorname{codim}(V(\mathfrak{a}),
  V(Ax_1+\cdots+Ax_r))=0.
\]
\end{proposition}

\begin{proof}
Soient $\mathfrak{p}_i$ $(1\leq i\leq m)$ les idéaux premiers minimaux de
l'anneau $A$, $\mathfrak{q}_j$ $(1\leq j\leq n)$ les idéaux premiers minimaux
parmi ceux qui contiennent $\mathfrak{a}$~; on a (14.2.1)
\[
  \operatorname{codim}(V(\mathfrak{a}),X)
  =\inf_j(\operatorname{codim}(V(\mathfrak{q}_j),X))
  =\inf_j(\dim(A_{\mathfrak{q}_j}))
\]
(16.1.3.2). L'hypothèse $r>0$ entraîne que $\mathfrak{a}$ n'est contenu dans
aucun des $\mathfrak{p}_i$~; par suite il existe $x\in\mathfrak{a}$ qui
n'appartient à aucun des $\mathfrak{p}_i$ (Bourbaki, \emph{Alg.~comm.},
chap.~II, \S~1, no~1, prop.~2). On a comme ci-dessus
\[
  \operatorname{codim}(V(\mathfrak{a}),V(Ax))
  =\inf_j(\dim((A/Ax)_{\mathfrak{q}_j/Ax}))
  =\inf_j(\dim(A_{\mathfrak{q}_j}/A_{\mathfrak{q}_j}x))~;
\]
mais comme $x/1$ n'est contenu dans aucun des idéaux premiers minimaux de
$A_{\mathfrak{q}_j}$, on a
$\dim(A_{\mathfrak{q}_j}/A_{\mathfrak{q}_j}x)
=\dim(A_{\mathfrak{q}_j})-1$ (16.3.4), d'où
$\operatorname{codim}(V(\mathfrak{a}),V(Ax))=r-1$. Pour la même raison, on a
$\dim(A/Ax)=\dim(A)-1$. Raisonnons maintenant par récurrence sur $r$~; si
$r>1$, il existe $r-1$ éléments $x'_i$ $(2\leq i\leq r)$ de $A'=A/Ax$,
faisant partie d'un système de paramètres de cet anneau, appartenant à
$\mathfrak{a}/Ax$ et tels que
\[
  \operatorname{codim}(V(\mathfrak{a}/Ax),
  V(A'x'_2+\cdots+A'x'_r))=0.
\]
Soit $(x'_i)_{r+1\leq i\leq s}$ un système d'éléments de $A'$ tel que
$(x'_i)_{2\leq i\leq s}$ soit un système de paramètres de $A'$~; pour tout
$i$ tel que $2\leq i\leq s$, soit $x_i$ un élément de la classe $x'_i$ dans
$A$, avec $x_i\in\mathfrak{a}$ pour $2\leq i\leq r$~; il est immédiat que
$A/(Ax+Ax_2+\cdots+Ax_s)$ est de longueur finie, et comme $\dim(A)=s$, on voit
que $x_1=x$ et les $x_i$ d'indices $i\geq2$ répondent aux conditions de
l'énoncé.
\end{proof}

\begin{proposition}[16.3.9]
\label{0.16.3.9-fr}
Soient $A$, $B$ deux anneaux locaux noethériens, $\mathfrak{m}$ l'idéal maximal
de $A$, $k$ son corps résiduel, $\varphi:A\to B$ un homomorphisme local.
\begin{enumerate}
  \item[\rm{(i)}] On a
  \begin{equation}
  \label{0.16.3.9.1-fr}
  \tag{16.3.9.1}
    \dim(B)\leq\dim(A)+\dim(B\otimes_A k).
  \end{equation}
  \item[\rm{(ii)}] Soient $M\neq0$ un $A$-module de type fini, $N\neq0$ un
  $B$-module de type fini. On a
  \begin{equation}
  \label{0.16.3.9.2-fr}
  \tag{16.3.9.2}
    \dim_B(M\otimes_A N)
    \leq\dim_A(M)+\dim_{B\otimes_A k}(N\otimes_A k).
  \end{equation}
\end{enumerate}
\end{proposition}

\begin{proof}
Soient $m$ la dimension de $A$, $(s_i)_{1\leq i\leq m}$ un système de
paramètres de $A$ (16.3.6), de sorte que si
$\mathfrak{a}=\sum_i s_iA$, $A/\mathfrak{a}$ est un $A$-module de longueur
finie. Alors l'anneau $A/\mathfrak{a}$ est artinien, donc l'idéal maximal
$\mathfrak{m}/\mathfrak{a}$ de cet anneau est nilpotent~; l'idéal
$\mathfrak{m}B/\mathfrak{a}B$ de $B/\mathfrak{a}B$, image de
$(\mathfrak{m}/\mathfrak{a})\otimes_A B$, est donc lui aussi nilpotent, et par
suite on a
$\dim(B/\mathfrak{a}B)=\dim(B/\mathfrak{m}B)$ (les spectres de ces deux
anneaux ayant même espace sous-jacent). D'autre part,
$B/\mathfrak{a}B=B/(s_1B+\cdots+s_mB)$, et les images des $s_i$ dans $B$
appartiennent à l'idéal maximal de $B$. Par suite (16.3.7), on a
\[
  \dim(B)\leq m+\dim(B\otimes_A k),
\]
ce qui n'est autre que (16.3.9.1).

Pour prouver (16.3.9.2), notons que si $\mathfrak{r}$ (resp.~$\mathfrak{s}$)
est l'annulateur de $M$ (resp.~$N$)
\oldpage[0\textsubscript{IV-1}]{32}
dans $A$ (resp.~$B$), on a
$\dim_A(M)=\dim(A/\mathfrak{r})$, $\dim_B(N)=\dim(B/\mathfrak{s})$ par
définition (16.1.7)~; d'autre part, $\operatorname{Supp}(M\otimes_A N)$ est un
sous-ensemble fermé de $\operatorname{Spec}(B)$ qui s'identifie (I, 9.1.13) à
\[
  \operatorname{Supp}(M)\times_{\operatorname{Spec}(A)}\operatorname{Supp}(N)
  =\operatorname{Spec}((B/\mathfrak{s})\otimes_A(A/\mathfrak{r}))
  =\operatorname{Spec}(B/(\mathfrak{s}+\mathfrak{r}B)).
\]
Comme $N\otimes_A k$ est un $B$-module de type fini, on a
$\dim_{B\otimes_A k}(N\otimes_A k)=\dim_B(N\otimes_A k)$ (16.1.9), et
$\operatorname{Supp}(N\otimes_A k)$ est égal à
$\operatorname{Spec}(B/(\mathfrak{s}+\mathfrak{m}B))$ par le même raisonnement
que ci-dessus. Appliquant (16.3.9.1) à $A/\mathfrak{r}$ et à
$B/(\mathfrak{s}+\mathfrak{r}B)$ il vient
\[
  \dim(B/(\mathfrak{s}+\mathfrak{r}B))
  \leq\dim(A/\mathfrak{r})+
  \dim(B/(\mathfrak{s}+\mathfrak{m}B)),
\]
ce qui n'est autre que (16.3.9.2).
\end{proof}

\begin{corollary}[16.3.10]
\label{0.16.3.10-fr}
Sous les hypothèses de (16.3.9), si $\mathfrak{m}B$ est un idéal de définition
de $B$, on a $\dim(B)\leq\dim(A)$~; si de plus $A$ est intègre et
$\dim(A)=\dim(B)$, $\varphi$ est injectif.
\end{corollary}

\begin{proof}
La première assertion résulte de (16.3.9) puisque alors
$\dim(B/\mathfrak{m}B)=0$. On peut d'ailleurs remplacer $A$ par
$\varphi(A)$, donc $\dim(B)\leq\dim(\varphi(A))$~; si
$\mathfrak{a}=\operatorname{Ker}(\varphi)\neq0$, et si $A$ est intègre, on a
$\dim(\varphi(A))=\dim(A/\mathfrak{a})<\dim(A)$ (16.1.2.2), donc on ne peut
alors avoir $\dim(A)=\dim(B)$ que si $\mathfrak{a}=0$.
\end{proof}

\subsection[Profondeur et coprofondeur]{Profondeur et
coprofondeur\protect\footnote{Dans la 3e partie du chap.~III, nous développons
la notion de profondeur du point de vue cohomologique et dans un cadre plus
général.}}
\label{subsection:0.16.4-fr}

\begin{proposition}[16.4.1]
\label{0.16.4.1-fr}
Soient $A$ un anneau local noethérien, $\mathfrak{m}$ son idéal maximal, $M$ un
$A$-module de type fini. Alors toute suite $M$-régulière
$(x_i)_{1\leq i\leq r}$ (15.1.7) formée d'éléments de $\mathfrak{m}$, fait
partie d'un système de paramètres pour $M$.
\end{proposition}

\begin{proof}
Raisonnons par récurrence sur $r$~; considérons le $A$-module
\[
  N=M/(x_1M+\cdots+x_{r-1}M)~;
\]
par hypothèse $x_1,\ldots,x_{r-1}$ font partie d'un système de paramètres pour
$M$, donc $\dim(N)=\dim(M)-(r-1)$ (16.3.7). Comme par hypothèse l'homothétie
de rapport $x_r$ dans $N$ est injective, $x_r$ n'appartient à aucun des idéaux
premiers associés à $N$, donc (16.3.4) on a
$\dim(N/x_rN)=\dim(N)-1$, ce qui s'écrit aussi
\[
  \dim(M/(x_1M+\cdots+x_rM))=\dim(M)-r~;
\]
la conclusion résulte donc de (16.3.7).
\end{proof}

Pour un $A$-module $M\neq0$, une suite $M$-régulière d'éléments de
$\mathfrak{m}$ a donc \emph{au plus} $\dim(M)$ éléments.

\begin{lemma}[16.4.2]
\label{0.16.4.2-fr}
Soient $A$ un anneau local noethérien, $\mathfrak{m}$ son idéal maximal,
$k=A/\mathfrak{m}$ son corps résiduel, $M$ un $A$-module. Pour qu'une suite
$M$-régulière $(x_i)_{1\leq i\leq n}$ d'éléments de $\mathfrak{m}$ soit
maximale, il faut et il suffit que l'on ait
\[
  \operatorname{Hom}_A(k,M/(x_1M+\cdots+x_nM))\neq0.
\]
\end{lemma}

\begin{proof}
Dire que $(x_i)$ est maximale signifie que, pour aucun $x\in\mathfrak{m}$,
l'homothétie de rapport $x$ dans $N=M/(x_1M+\cdots+x_nM)$ n'est injective, ou
encore que $\mathfrak{m}$ est contenu dans la réunion des idéaux premiers
associés à $N$ (Bourbaki, \emph{Alg.~comm.}, chap.~IV, \S~1, no~1, cor.~2 de
la prop.~2), donc égal à un de ces derniers, puisque c'est un idéal
\oldpage[0\textsubscript{IV-1}]{33}
maximal (Bourbaki, \emph{Alg.~comm.}, chap.~II, \S~1, no~1, prop.~2)~;
autrement dit, il y a un élément $z\in N$ dont l'annulateur est
$\mathfrak{m}$, et par suite le sous-module $Az$ de $N$ est isomorphe à
$A/\mathfrak{m}=k$~; d'où le lemme.
\end{proof}

\begin{lemma}[16.4.3]
\label{0.16.4.3-fr}
Soient $A$ un anneau local noethérien, $\mathfrak{m}$ son idéal maximal,
$k=A/\mathfrak{m}$ son corps résiduel, $M$ un $A$-module,
$(x_i)_{1\leq i\leq n}$ une suite $M$-régulière d'éléments de
$\mathfrak{m}$. Alors les $A$-modules
\[
  \operatorname{Hom}_A(k,M/(x_1M+\cdots+x_nM))
  \quad\text{et}\quad \operatorname{Ext}^n_A(k,M)
\]
sont isomorphes, et pour $n\geq1$, ils sont aussi isomorphes à
$\operatorname{Ext}^{n-1}_A(k,M/x_1M)$.
\end{lemma}

\begin{proof}
Raisonnons par récurrence sur $n$, la proposition étant évidente pour $n=0$.
Posons $N=M/x_1M$~; la suite $(x_i)_{2\leq i\leq n}$ est $N$-régulière, donc
\[
  \operatorname{Hom}_A(k,M/(x_1M+\cdots+x_nM))
  =\operatorname{Hom}_A(k,N/(x_2N+\cdots+x_nN))
\]
est isomorphe à $\operatorname{Ext}^{n-1}_A(k,N)$ en vertu de l'hypothèse de
récurrence. Considérons alors la suite exacte
$0\to M\xrightarrow{x_1}M\to M/x_1M=N\to0$~; on en déduit la suite exacte des
Ext
\begin{equation}
\label{0.16.4.3.1-fr}
\tag{16.4.3.1}
  \operatorname{Ext}^{n-1}_A(k,M)
  \to\operatorname{Ext}^{n-1}_A(k,N)
  \to\operatorname{Ext}^n_A(k,M)
  \xrightarrow{x_1}\operatorname{Ext}^n_A(k,M).
\end{equation}
En vertu de l'hypothèse de récurrence, $\operatorname{Ext}^{n-1}_A(k,M)$ est
isomorphe à
\[
  \operatorname{Hom}_A(k,M/(x_1M+\cdots+x_{n-1}M)),
\]
qui est nul, puisque $(x_i)_{1\leq i\leq n-1}$ n'est pas une suite
$M$-régulière maximale (16.4.2). D'autre part, comme
$x_1\in\mathfrak{m}$, l'homothétie de rapport $x_1$ dans le $A$-module
$\operatorname{Hom}_A(k,T)$ est nulle pour tout $A$-module $T$, donc il en est
de même de l'homothétie de rapport $x_1$ dans tout $A$-module
$\operatorname{Ext}^n_A(k,T)$~; les assertions du lemme découlent donc
aussitôt de la suite exacte (16.4.3.1).
\end{proof}

\begin{corollary}[16.4.4]
\label{0.16.4.4-fr}
Soient $A$ un anneau local noethérien, $\mathfrak{m}$ son idéal maximal,
$k=A/\mathfrak{m}$ son corps résiduel, $M$ un $A$-module de type fini
$\neq0$. Toutes les suites $M$-régulières maximales d'éléments de
$\mathfrak{m}$ ont le même nombre d'éléments, qui est le plus petit entier $n$
tel que $\operatorname{Ext}^n_A(k,M)\neq0$.
\end{corollary}

\begin{definition}[16.4.5]
\label{0.16.4.5-fr}
Soient $A$ un anneau local noethérien, $\mathfrak{m}$ son idéal maximal, $M$ un
$A$-module de type fini. On appelle \emph{profondeur} de $M$, et l'on note
$\operatorname{prof}_A M$ ou $\operatorname{prof}(M)$ la borne supérieure du
nombre d'éléments d'une suite $M$-régulière formée d'éléments de
$\mathfrak{m}$. On a donc $\operatorname{prof}(M)=+\infty$ si et seulement si
$M=0$, et, par (16.4.1)
\begin{equation}
\label{0.16.4.5.1-fr}
\tag{16.4.5.1}
  \operatorname{prof}(M)\leq\dim(M)\quad\text{si}\quad M\neq0.
\end{equation}
On a évidemment $\operatorname{prof}(M^k)=\operatorname{prof}(M)$ pour tout
$k>0$.
\end{definition}

\begin{proposition}[16.4.6]
\label{0.16.4.6-fr}
Soient $A$ un anneau local noethérien, $\mathfrak{m}$ son idéal maximal, $M$ un
$A$-module de type fini.
\begin{enumerate}
  \item[\rm{(i)}] Pour que $\operatorname{prof}(M)=0$, il faut et il suffit que
  $\mathfrak{m}\in\operatorname{Ass}(M)$ (ou encore qu'il existe un sous-module
  de $M$ isomorphe au corps résiduel $k=A/\mathfrak{m}$ de $A$).
  \item[\rm{(ii)}] Pour tout élément $M$-régulier $x$ appartenant à
  $\mathfrak{m}$, on a
  \begin{equation}
  \label{0.16.4.6.1-fr}
  \tag{16.4.6.1}
    \operatorname{prof}(M/xM)=\operatorname{prof}(M)-1.
  \end{equation}
  \item[\rm{(iii)}] Si $M\neq0$, on a
  \begin{equation}
  \label{0.16.4.6.2-fr}
  \tag{16.4.6.2}
    \operatorname{prof}(M)
    \leq\inf_{\mathfrak{p}\in\operatorname{Ass}(M)}\dim(A/\mathfrak{p})
    \leq\dim(M).
  \end{equation}
  \oldpage[0\textsubscript{IV-1}]{34}
  \item[\rm{(iv)}] On a
  $\operatorname{prof}_A(M)=
  \operatorname{prof}_{\widehat{A}}(\widehat{M})$.
\end{enumerate}
\end{proposition}

\begin{proof}
(i) Dire que $\operatorname{prof}(M)=0$ signifie qu'il n'existe aucun
$x\in\mathfrak{m}$ tel que l'homothétie de rapport $x$ dans $M$ soit
injective, autrement dit que tout $x\in\mathfrak{m}$ appartient à un idéal
premier associé à $M$ (Bourbaki, \emph{Alg.~comm.}, chap.~IV, \S~1, no~1,
cor.~2 de la prop.~2)~; mais $\mathfrak{m}$ ne peut être réunion des idéaux
premiers associés à $M$ que s'il est égal à l'un d'eux (Bourbaki,
\emph{Alg.~comm.}, chap.~II, \S~1, no~1, prop.~2).

(ii) Si $(x_i)_{1\leq i\leq m}$ est une suite $(M/xM)$-régulière formée
d'éléments de $\mathfrak{m}$, la suite formée de $x$ et des $x_i$ est
$M$-régulière, et il résulte de la définition qu'elle est maximale lorsque la
suite $(x_i)_{1\leq i\leq m}$ est une suite $(M/xM)$-régulière maximale, d'où
la conclusion.

(iii) Raisonnons par récurrence sur $n=\operatorname{prof}(M)$, l'assertion
étant évidente pour $n=0$. Nous utiliserons le lemme suivant~:
\end{proof}

\begin{lemma}[16.4.6.3]
\label{0.16.4.6.3-fr}
Soient $A$ un anneau noethérien, $M$ un $A$-module de type fini, $t\in A$ un
élément $M$-régulier. Alors, pour tout
$\mathfrak{p}\in\operatorname{Ass}(M)$, tout idéal premier minimal parmi ceux
qui contiennent $\mathfrak{p}+At$ est associé à $M/tM$.
\end{lemma}

\begin{proof}
On sait qu'il existe une suite exacte
\[
  0\to M'\to M\to M''\to0
\]
telle que $\operatorname{Ass}(M')=\{\mathfrak{p}\}$,
$\operatorname{Ass}(M'')=\operatorname{Ass}(M)-\{\mathfrak{p}\}$ (Bourbaki,
\emph{Alg.~comm.}, chap.~IV, \S~1, no~1, prop.~4)~; comme $t$ est
$M$-régulier, il est aussi $M'$-régulier et $M''$-régulier
(\emph{loc.~cit.}, \S~1, no~1, cor.~2 de la prop.~2). On en déduit que la
suite
\[
  0\to M'/tM'\to M/tM\to M''/tM''\to0
\]
est exacte (15.1.18). Par suite, on a
$\operatorname{Ass}(M'/tM')\subset\operatorname{Ass}(M/tM)$. Mais les idéaux
premiers de $A$ contenant $\mathfrak{p}$ sont les points de
$\operatorname{Supp}(M')$ (\emph{loc.~cit.}, \S~1, no~3, prop.~7), donc les
idéaux premiers contenant $\mathfrak{p}+At$ sont les points du support de
$M'/tM'=M'\otimes_A(A/tA)$ ($0_{\mathrm I}$, 1.7.5)~; un élément minimal de
l'ensemble de ces idéaux premiers appartient donc à
$\operatorname{Ass}(M'/tM')$ (Bourbaki, \emph{loc.~cit.}, \S~1, no~3,
cor.~1 de la prop.~7), et \emph{a fortiori} à
$\operatorname{Ass}(M/tM)$.
\end{proof}

\begin{proof}[Suite de la démonstration de (16.4.6)]
Ce lemme étant établi, revenons à la démonstration de (iii). Soit
$x\in\mathfrak{m}$ un élément $M$-régulier, de sorte que
$\operatorname{prof}(M/xM)=n-1$ par (ii)~; pour tout
$\mathfrak{p}\in\operatorname{Ass}(M)$, il résulte du lemme qu'il y a un idéal
$\mathfrak{p}'$ associé à $M/xM$ et contenant $\mathfrak{p}+xA$~; mais comme
$x$ est $M$-régulier, on a $x\notin\mathfrak{p}$, d'où
$\mathfrak{p}'\neq\mathfrak{p}$ et par suite
$\dim(A/\mathfrak{p}')\leq\dim(A/\mathfrak{p})-1$. Or l'hypothèse de
récurrence montre que $n-1\leq\dim(A/\mathfrak{p}')$~; on en conclut que
$n\leq\dim(A/\mathfrak{p})$ pour tout
$\mathfrak{p}\in\operatorname{Ass}(M)$.

(iv) On peut se borner au cas où $M\neq0$. Rappelons que
$\widehat{M}=M\otimes_A\widehat{A}$ à un isomorphisme canonique près
($0_{\mathrm I}$, 7.3.3)~; comme $\widehat{A}$ est un $A$-module plat, on a des
isomorphismes canoniques
\[
  \operatorname{Ext}^i_A(k,M)\otimes_A\widehat{A}
  \xrightarrow{\sim}
  \operatorname{Ext}^i_{\widehat{A}}
  (k,M\otimes_A\widehat{A})
\]
puisque $k=k\otimes_A\widehat{A}$ à un isomorphisme canonique près
($0_{\mathrm{III}}$, 12.3.5). L'assertion (iv) résulte donc de (16.4.4) et du
fait que $\widehat{A}$ est un $A$-module fidèlement plat
($0_{\mathrm I}$, 7.3.5).
\end{proof}

\begin{env}[16.4.6.4]
\label{0.16.4.6.4-fr}
On notera que si $\mathfrak{p}$ est un idéal premier de $A$, on peut avoir
aussi bien
$\operatorname{prof}_{A_{\mathfrak{p}}}(M_{\mathfrak{p}})
<\operatorname{prof}_A(M)$ que
$\operatorname{prof}_{A_{\mathfrak{p}}}(M_{\mathfrak{p}})
>\operatorname{prof}_A(M)$~; pour un exemple du premier cas,
\oldpage[0\textsubscript{IV-1}]{35}
il suffit de prendre pour $A$ un anneau local intègre de dimension $\geq1$~; on a
$\operatorname{prof}(A)\geq1$ mais $\operatorname{prof}(K)=0$ pour le corps
des fractions $K$ de $A$. Pour un exemple du second cas, considérons un
anneau local noethérien intègre $A_0$ de dimension $\geq2$, d'idéal maximal
$\mathfrak{m}_0$ et soit $k=A_0/\mathfrak{m}_0$ son corps résiduel~; soit
$A=A_0\oplus\mathfrak{J}$ l'extension triviale de $A_0$ par le $A_0$-module
$k$ (18.2.3), l'idéal $\mathfrak{J}$ étant $A_0$-isomorphe à $k$ (de sorte
que la multiplication dans $A$ est donnée par
$(x,z)(x',z')=(xx',xz'+x'z)$). Il est clair que $\mathfrak{J}$ est le
nilradical de $A$ et $A/\mathfrak{J}$ est isomorphe à $A_0$, de sorte que
tout idéal premier $\mathfrak{p}$ de $A$ est de la forme
$\mathfrak{p}_0\oplus\mathfrak{J}$, où $\mathfrak{p}_0$ est un idéal premier
de $A_0$~; en particulier $\mathfrak{m}=\mathfrak{m}_0\oplus\mathfrak{J}$
est l'unique idéal maximal de $A$. Comme tout élément de $\mathfrak{m}_0$
annule $\mathfrak{J}$, on a
$\operatorname{prof}(A)=\operatorname{prof}(A_{\mathfrak{m}})=0$~; au
contraire, si $\mathfrak{p}_0\neq\mathfrak{m}_0$, les éléments de
$\mathfrak{J}$ sont annulés par ceux de
$\mathfrak{m}_0-\mathfrak{p}_0$, donc $A_{\mathfrak{p}}$ est canoniquement
isomorphe à $(A_0)_{\mathfrak{p}_0}$, et comme $A_0$ est intègre,
$\operatorname{prof}(A_{\mathfrak{p}})
=\operatorname{prof}((A_0)_{\mathfrak{p}_0})\geq1$ si
$\mathfrak{p}_0\neq0$.
\end{env}

\begin{corollary}[16.4.7]
\label{0.16.4.7-fr}
Pour qu'un anneau local noethérien réduit $A$ soit de profondeur $0$, il faut
et il suffit que $A$ soit un corps.
\end{corollary}

\begin{proof}
En effet, comme l'intersection des idéaux premiers minimaux de $A$ est $0$,
$A$ n'a pas d'idéaux premiers associés immergés~; en vertu de (16.4.6, (i)),
si $\operatorname{prof}(A)=0$, l'idéal maximal $\mathfrak{m}$ de $A$ est
aussi idéal premier minimal, donc c'est le seul idéal premier de $A$, et
comme $A$ est réduit, $\mathfrak{m}=0$.
\end{proof}

\begin{proposition}[16.4.8]
\label{0.16.4.8-fr}
Soient $A$, $B$ deux anneaux locaux noethériens,
$\rho:A\to B$ un homomorphisme local, $M$ un $B$-module tel que
$M_{[\rho]}$ soit un $A$-module de type fini. Alors on a
\begin{equation}
\label{0.16.4.8.1-fr}
\tag{16.4.8.1}
  \operatorname{prof}_A(M_{[\rho]})=\operatorname{prof}_B(M).
\end{equation}
\end{proposition}

\begin{proof}
On peut se borner au cas $M\neq0$~; supposons que
$\operatorname{prof}_A(M_{[\rho]})=n$. Si $(x_i)_{1\leq i\leq n}$ est une
suite $(M_{[\rho]})$-régulière maximale formée d'éléments de l'idéal maximal
de $A$, les $\rho(x_i)$ constituent une suite $M$-régulière formée d'éléments
de l'idéal maximal de $B$, et on a donc, en posant
$N=M/(\rho(x_1)M+\cdots+\rho(x_n)M)$,
$\operatorname{prof}_B(N)=\operatorname{prof}_B(M)-n$ (16.4.6, (ii))~; de
même
$N_{[\rho]}=M_{[\rho]}/(x_1M_{[\rho]}+\cdots+x_nM_{[\rho]})$ et
$\operatorname{prof}_A(N_{[\rho]})=0$~; on est ainsi ramené à démontrer la
proposition lorsque $n=0$. Soit $k$ le corps résiduel de $A$~; comme $k$ est
un quotient de $A$, le $A$-module
$P=\operatorname{Hom}_A(k,M_{[\rho]})$ est un sous-module de
$\operatorname{Hom}_A(A,M_{[\rho]})=M_{[\rho]}$ et l'hypothèse entraîne que
$P\neq0$ (16.4.6, (i))~; en outre, $P$ est un $A$-module de type fini
(puisque $M_{[\rho]}$ est un $A$-module de type fini), donc un $k$-espace
vectoriel de type fini, étant annulé par l'idéal maximal de $A$, et finalement
un $A$-module de longueur finie. D'autre part, $P$ est aussi un sous-module du
$B$-module $M$, et puisqu'il est de longueur finie en tant que $A$-module, il
l'est \emph{a fortiori} en tant que $B$-module. Comme il est $\neq0$, il
contient un sous-$B$-module simple, c'est-à-dire isomorphe au corps résiduel de
$B$~; on conclut alors de (16.4.6, (i)) que l'on a bien
$\operatorname{prof}_B(M)=0$.
\end{proof}

\begin{definition}[16.4.9]
\label{0.16.4.9-fr}
Soient $A$ un anneau local noethérien, $M$ un $A$-module de type fini. Si
$M\neq0$, on appelle \emph{coprofondeur} de $M$ et on note
$\operatorname{coprof}_A(M)$ ou $\operatorname{coprof}(M)$, l'entier fini
$\dim_A(M)-\operatorname{prof}_A(M)\geq0$ (16.4.5.1). Si $M=0$, on pose
$\operatorname{coprof}(M)=0$.

On a $\operatorname{coprof}(M^k)=\operatorname{coprof}(M)$ pour tout $k>0$.
\end{definition}

\begin{proposition}[16.4.10]
\label{0.16.4.10-fr}
Soient $A$ un anneau local noethérien, $M$ un $A$-module de type fini.
\oldpage[0\textsubscript{IV-1}]{36}
\begin{enumerate}
  \item[(i)] Pour tout élément $M$-régulier $x$ appartenant à l'idéal maximal
  de $A$, on a
  $\operatorname{coprof}(M/xM)=\operatorname{coprof}(M)$.
  \item[(ii)] On a
  $\operatorname{coprof}_A(M)=\operatorname{coprof}_{\widehat A}(\widehat M)$.
\end{enumerate}
\end{proposition}

\begin{proof}
En effet, (i) résulte de (16.3.4) et (16.4.6, (ii)), et (ii) de (16.2.4) et
(16.4.6, (iv)).
\end{proof}

Nous montrerons plus tard (IV, 6.11.5) que pour tout idéal premier
$\mathfrak{p}$ de $A$, on a
$\operatorname{coprof}_{A_{\mathfrak{p}}}(M_{\mathfrak{p}})
\leq\operatorname{coprof}_A(M)$.

\begin{proposition}[16.4.11]
\label{0.16.4.11-fr}
Sous les hypothèses de (16.4.8), on a
\begin{equation}
\label{0.16.4.11.1-fr}
\tag{16.4.11.1}
  \operatorname{coprof}_A(M_{[\rho]})=\operatorname{coprof}_B(M).
\end{equation}
\end{proposition}

\begin{proof}
Cela résulte de (16.1.9) et (16.4.8).
\end{proof}

\subsection{Modules de Cohen-Macaulay}
\label{subsection:0.16.5-fr}

\begin{definition}[16.5.1]
\label{0.16.5.1-fr}
Soit $A$ un anneau local noethérien. On dit qu'un $A$-module de type fini $M$
est un \emph{module de Cohen-Macaulay} si
$\operatorname{coprof}(M)=0$, autrement dit si $M=0$ ou si $M\neq0$ et
$\dim(M)=\operatorname{prof}(M)$. On dit que $A$ est un
\emph{anneau de Cohen-Macaulay} si le $A$-module $A$ est un module de
Cohen-Macaulay.
\end{definition}

Un $A$-module de longueur finie est un module de Cohen-Macaulay puisqu'il est
de dimension $0$ (16.1.10). Un anneau local réduit $A$ de dimension $1$ est un
anneau de Cohen-Macaulay, car alors son idéal maximal $\mathfrak{m}$ ne peut
être associé à $A$ (Bourbaki, \emph{Alg.~comm.}, chap.~IV, \S~2, no~5,
prop.~10), donc $\operatorname{prof}(A)=1$ en vertu de (16.4.6, (i)). Si $A$
est un \emph{anneau local intègre et intégralement clos de dimension $\geq2$},
on a $\operatorname{prof}(A)\geq2$~: en effet, si $x\neq0$ est un élément de
$\mathfrak{m}$, on sait (Bourbaki, \emph{Alg.~comm.}, chap.~VII, \S~1, no~4,
prop.~8) que les idéaux premiers associés à $A/xA$ sont de hauteur $1$, donc
distincts de $\mathfrak{m}$, et leur réunion est par suite distincte de
$\mathfrak{m}$ (Bourbaki, \emph{Alg.~comm.}, chap.~II, \S~1, no~1, prop.~2)~;
il existe donc des suites $(x,y)$ qui sont $A$-régulières et formées
d'éléments de $\mathfrak{m}$, d'où notre assertion. On en conclut qu'un anneau
local intégralement clos de dimension $2$ est un anneau de Cohen-Macaulay.

\begin{proposition}[16.5.2]
\label{0.16.5.2-fr}
Soient $A$ un anneau local noethérien, $M$ un $A$-module de type fini. Pour que
$M$ soit un module de Cohen-Macaulay, il faut et il suffit que $\widehat M$
soit un $\widehat A$-module de Cohen-Macaulay.
\end{proposition}

\begin{proof}
Cela résulte aussitôt de (16.4.10, (ii)).
\end{proof}

\begin{proposition}[16.5.3]
\label{0.16.5.3-fr}
Sous les hypothèses de (16.4.8), pour que $M$ soit un $B$-module de
Cohen-Macaulay, il faut et il suffit que $M_{[\rho]}$ soit un $A$-module de
Cohen-Macaulay.
\end{proposition}

\begin{proof}
Cela résulte de (16.4.11).
\end{proof}

\begin{proposition}[16.5.4]
\label{0.16.5.4-fr}
Soient $A$ un anneau local noethérien, $M$ un $A$-module de type fini $\neq0$.
Si $M$ est un $A$-module de Cohen-Macaulay, on a
$\dim(A/\mathfrak{p})=\dim(M)$ pour tout idéal premier
$\mathfrak{p}\in\operatorname{Ass}(M)$~; en particulier aucun idéal premier
associé à $M$ n'est immergé.
\end{proposition}

\begin{proof}
Cela résulte aussitôt des inégalités (16.4.6.2).
\end{proof}

On peut encore dire que $M$ (ou $\operatorname{Supp}(M)$) est
\emph{équidimensionnel} (16.1.7).

\oldpage[0\textsubscript{IV-1}]{37}

\begin{proposition}[16.5.5]
\label{0.16.5.5-fr}
Soient $A$ un anneau local noethérien, $\mathfrak{m}$ son idéal maximal, $M$
un $A$-module de type fini $\neq0$, $x$ un élément de $\mathfrak{m}$.
Supposons que $M$ soit un module de Cohen-Macaulay~; alors les conditions
suivantes sont équivalentes~:
\begin{enumerate}
  \item[a)] $x$ est $M$-régulier~;
  \item[b)] $\dim(M/xM)=\dim(M)-1$~;
  \item[c)] $x$ n'appartient à aucun élément minimal de
  $\operatorname{Supp}(M)$.
\end{enumerate}
En outre, $M/xM$ est alors un module de Cohen-Macaulay.
\end{proposition}

\begin{proof}
On a vu (16.5.4) que tous les éléments $\mathfrak{p}$ de
$\operatorname{Ass}(M)$ sont minimaux dans $\operatorname{Supp}(M)$ et que
$\dim(A/\mathfrak{p})=\dim(M)$~; l'équivalence de a) et c) résulte donc de
Bourbaki, \emph{Alg.~comm.}, chap.~IV, \S~1, no~1, cor.~2 de la prop.~2~;
l'équivalence de b) et c) résulte de (16.3.4). Le fait que $M/xM$ soit un
module de Cohen-Macaulay résulte enfin de (16.4.10, (i)).
\end{proof}

\begin{corollary}[16.5.6]
\label{0.16.5.6-fr}
Sous les hypothèses de (16.5.5), soit $(x_i)_{1\leq i\leq r}$ une suite
d'éléments de $\mathfrak{m}$. Les conditions suivantes sont équivalentes~:
\begin{enumerate}
  \item[a)] La suite $(x_i)$ est $M$-régulière.
  \item[b)] On a
  $\dim(M/(x_1M+\cdots+x_rM))=\dim(M)-r$.
  \item[c)] Les $x_i$ font partie d'un système de paramètres de $M$.
\end{enumerate}
En outre, lorsque ces conditions sont satisfaites,
$N=M/(x_1M+\cdots+x_rM)$ est un module de Cohen-Macaulay, donc, pour tout
idéal $\mathfrak{p}\in\operatorname{Ass}(N)$, on a
$\dim(A/\mathfrak{p})=\dim(M)-r$.
\end{corollary}

\begin{proof}
On sait déjà que b) et c) sont équivalentes (16.3.6) et que a) entraîne c)
(16.4.1) sans hypothèse sur $M$. Pour voir que c) implique a) et prouver la
dernière assertion de l'énoncé, raisonnons par récurrence sur $r$, l'assertion
étant triviale pour $r=0$. Puisque $x_1$ fait partie d'un système de
paramètres, $\dim(M/x_1M)=\dim(M)-1$ (16.3.6), donc $x_1$ est $M$-régulier et
$M/x_1M$ est un module de Cohen-Macaulay par (16.5.5). Comme la suite
$(x_i)_{2\leq i\leq r}$ fait partie d'un système de paramètres pour
$P=M/x_1M$, l'hypothèse de récurrence montre que $(x_i)_{2\leq i\leq r}$ est
une suite $P$-régulière et que
$P/(x_2P+\cdots+x_rP)=N$ est un module de Cohen-Macaulay~; en outre
$(x_i)_{1\leq i\leq r}$ est $M$-régulière.
\end{proof}

\begin{remark}[16.5.7]
\label{0.16.5.7-fr}
Il résulte de (16.5.6) qu'il y a identité entre les \emph{systèmes de
paramètres pour $M$} et les \emph{suites $M$-régulières maximales d'éléments de
$\mathfrak{m}$} lorsque $M$ est un module de Cohen-Macaulay. Inversement, il
résulte de (16.4.1) que si un $A$-module de type fini non nul est tel qu'il y
ait une suite $M$-régulière qui soit un système de paramètres pour $M$, alors
$M$ est un module de Cohen-Macaulay. La dernière propriété énoncée dans
(16.5.6) \emph{caractérise aussi} les modules de Cohen-Macaulay~:
\end{remark}

\begin{proposition}[16.5.8]
\label{0.16.5.8-fr}
Soient $A$ un anneau local noethérien, $M$ un $A$-module de type fini. On
suppose que pour toute suite $(x_i)_{1\leq i\leq r}$ d'éléments de $A$
faisant partie d'un système de paramètres pour $M$, et tout idéal premier
$\mathfrak{p}$ associé à
$N=M/(x_1M+\cdots+x_rM)$, on ait
$\dim(A/\mathfrak{p})=\dim(N)$. Alors $M$ est un module de Cohen-Macaulay.
\end{proposition}

\begin{proof}
Soit $(y_i)_{1\leq i\leq n}$ un système de paramètres pour $M$~; il suffit de
montrer que la suite $(y_i)$ est $M$-régulière. Raisonnons par récurrence sur
$n=\dim(M)$, la proposition étant triviale pour $n=0$. Par hypothèse
(appliquée avec $r=0$) les idéaux premiers $\mathfrak{p}_j$ associés à $M$
sont tous tels que $\dim(A/\mathfrak{p}_j)=\dim(M)$~; comme $y_1$ fait partie
d'un système de paramètres
\oldpage[0\textsubscript{IV-1}]{38}
pour $M$, il résulte de (16.3.7) et (16.3.4) que $y_1$ n'appartient à aucun
des $\mathfrak{p}_j$, donc est $M$-régulier. Si on pose $P=M/y_1M$,
$(y_i)_{2\leq i\leq n}$ est un système de paramètres pour $P$, et il est
immédiat que $P$ vérifie l'hypothèse de l'énoncé~; appliquant l'hypothèse de
récurrence, on voit donc que $(y_i)_{2\leq i\leq n}$ est une suite
$P$-régulière, ce qui démontre la proposition.
\end{proof}

\begin{proposition}[16.5.9]
\label{0.16.5.9-fr}
Soient $A$ un anneau local noethérien, $M$ un $A$-module de type fini~; on
suppose que $M$ soit un module de Cohen-Macaulay. Soit
$\mathfrak{p}\in\operatorname{Supp}(M)$, et posons
$r=\dim(M)-\dim(M/\mathfrak{p}M)$. Alors il existe une suite $M$-régulière
$(x_i)_{1\leq i\leq r}$ formée d'éléments de $\mathfrak{p}$~; pour toute
suite ayant ces propriétés, on a
\[
  \dim(M/(x_1M+\cdots+x_rM))
  =\dim(M/\mathfrak{p}M)=\dim(A/\mathfrak{p}),
\]
et $\mathfrak{p}$ est un élément minimal de
$\operatorname{Ass}(M/(x_1M+\cdots+x_rM))$.
\end{proposition}

\begin{proof}
Prouvons d'abord l'existence des $x_i$~; il n'y a rien à démontrer pour
$r=0$. Pour $r>0$, procédons par récurrence sur $r$. Comme pour tout idéal
premier $\mathfrak{q}_i$ associé à $M$ on a
$\dim(A/\mathfrak{q}_i)=\dim(M)$ (16.5.4), l'hypothèse
$\dim(M/\mathfrak{p}M)<\dim(M)$ entraîne que $\mathfrak{p}$ ne peut être
contenu dans un des $\mathfrak{q}_i$, ces derniers étant les idéaux premiers
minimaux parmi ceux qui contiennent $\operatorname{Ann}(M)$ (16.1.2.2 et
16.1.7)~; on en conclut que $\mathfrak{p}$ n'est pas non plus contenu dans la
réunion des $\mathfrak{q}_i$ (Bourbaki, \emph{Alg.~comm.}, chap.~II, \S~1,
no~1, prop.~2), et par suite il existe un $x_1\in\mathfrak{p}$ qui est
$M$-régulier (Bourbaki, \emph{Alg.~comm.}, chap.~IV, \S~1, no~1, cor.~2 de
la prop.~2). On déduit donc de (16.5.5) que $N=M/x_1M$ est un module de
Cohen-Macaulay, de dimension $\dim(M)-1$~; comme $x_1\in\mathfrak{p}$, on a
d'ailleurs $N/\mathfrak{p}N=M/\mathfrak{p}M$, d'où
$\dim(N)-\dim(N/\mathfrak{p}N)=r-1$~; on peut par suite appliquer à $N$
l'hypothèse de récurrence, et si $(x_i)_{2\leq i\leq r}$ est une suite
$N$-régulière formée d'éléments de $\mathfrak{p}$, il est clair que
$(x_i)_{1\leq i\leq r}$ est une suite $M$-régulière formée d'éléments de
$\mathfrak{p}$.

Si on pose $P=M/(x_1M+\cdots+x_rM)$, on a
$P/\mathfrak{p}P=M/\mathfrak{p}M$ puisque les $x_i$ sont dans
$\mathfrak{p}$, et $\dim(P)=\dim(P/\mathfrak{p}P)$ par (16.5.6)~; mais comme
$P$ est un module de Cohen-Macaulay, on a aussi
$\dim(P)=\dim(A/\mathfrak{p}')$ pour tout
$\mathfrak{p}'\in\operatorname{Ass}(P)$ (16.5.4). Comme
$\mathfrak{p}\in\operatorname{Supp}(P)$ ($0_{\mathrm I}$, 1.7.5),
$\mathfrak{p}$ contient un des idéaux
$\mathfrak{p}'\in\operatorname{Ass}(P)$ (Bourbaki, \emph{Alg.~comm.},
chap.~IV, \S~1, no~4, th.~2), et comme
\[
  \dim(P)=\dim(P/\mathfrak{p}P)
  \leq\dim(A/\mathfrak{p})
  \leq\dim(A/\mathfrak{p}')=\dim(P),
\]
on a nécessairement $\mathfrak{p}=\mathfrak{p}'$ (16.1.2.2), ce qui achève
la démonstration.
\end{proof}

\begin{corollary}[16.5.10]
\label{0.16.5.10-fr}
Sous les mêmes hypothèses que dans (16.5.9)~:
\begin{enumerate}
  \item[(i)] $M_{\mathfrak{p}}$ est un $A_{\mathfrak{p}}$-module de
  Cohen-Macaulay.
  \item[(ii)] On a
  \begin{equation}
  \label{0.16.5.10.1-fr}
  \tag{16.5.10.1}
    \dim(M)=\dim_A(M/\mathfrak{p}M)
    +\dim_{A_{\mathfrak{p}}}(M_{\mathfrak{p}}).
  \end{equation}
\end{enumerate}
\end{corollary}

\begin{proof}
Avec les notations de (16.5.9), soient $x_i'$ les images canoniques des $x_i$
dans $A_{\mathfrak{p}}$ $(1\leq i\leq r)$~; comme les $x_i'$ appartiennent à
l'idéal maximal de $A_{\mathfrak{p}}$ et forment une suite
$M_{\mathfrak{p}}$-régulière par platitude (15.1.14), on a
\[
  \operatorname{prof}_{A_{\mathfrak{p}}}(M_{\mathfrak{p}})
  \geq r=\dim(M)-\dim(M/\mathfrak{p}M)
  \geq\dim_{A_{\mathfrak{p}}}(M_{\mathfrak{p}})
\]
(en vertu de (16.1.8.1)). Mais compte tenu de (16.4.5.1), les trois termes de
cette inégalité sont nécessairement égaux, d'où le corollaire.
\end{proof}

\oldpage[0\textsubscript{IV-1}]{39}

\begin{corollary}[16.5.11]
\label{0.16.5.11-fr}
Soit $A$ un anneau local noethérien~; supposons qu'il existe un $A$-module de
type fini $M$ de support $\operatorname{Spec}(A)$, qui soit un module de
Cohen-Macaulay. Alors, pour tout idéal premier $\mathfrak{p}$ de $A$ on a
\begin{equation}
\label{0.16.5.11.1-fr}
\tag{16.5.11.1}
  \dim(A)=\dim(A/\mathfrak{p})+\dim(A_{\mathfrak{p}}).
\end{equation}
\end{corollary}

\begin{proof}
En effet, $\operatorname{Supp}(M/\mathfrak{p}M)
=\operatorname{Spec}(A/\mathfrak{p})$ ($0_{\mathrm I}$, 1.7.5) et
$\operatorname{Supp}(M_{\mathfrak{p}})
=\operatorname{Spec}(A_{\mathfrak{p}})$~; la relation (16.5.11.1) est donc un
cas particulier de (16.5.10.1).
\end{proof}

\begin{corollary}[16.5.12]
\label{0.16.5.12-fr}
Tout anneau quotient d'un anneau local noethérien $A$ vérifiant les conditions
de (16.5.11) (en particulier tout anneau quotient d'un anneau local de
Cohen-Macaulay) est \emph{caténaire}.
\end{corollary}

\begin{proof}
Il suffit évidemment de prouver que $A$ lui-même est caténaire c'est-à-dire
que, pour deux idéaux premiers $\mathfrak{p}\subset\mathfrak{q}$ de $A$, on a
(16.1.4.2)
\[
  \dim(A_{\mathfrak{q}})
  =\dim(A_{\mathfrak{p}})
  +\dim(A_{\mathfrak{q}}/\mathfrak{p}A_{\mathfrak{q}}).
\]
Or, en vertu de (16.5.10, (i)) $A_{\mathfrak{q}}$ vérifie les mêmes
hypothèses que $A$, et la relation précédente n'est donc autre que
(16.5.11.1) appliquée à $A_{\mathfrak{q}}$ et à l'idéal premier
$\mathfrak{p}A_{\mathfrak{q}}$ de $A_{\mathfrak{q}}$.
\end{proof}

\begin{env}[16.5.13]
\label{0.16.5.13-fr}
Soient $A$ un anneau noethérien, $M$ un $A$-module de type fini. On dit que
$M$ est un \emph{$A$-module de Cohen-Macaulay} si, pour tout idéal premier
$\mathfrak{p}$ de $A$, $M_{\mathfrak{p}}$ est un $A_{\mathfrak{p}}$-module
de Cohen-Macaulay~; en vertu de (16.5.10) cette définition coïncide avec
(16.5.1) lorsque $A$ est \emph{local}. On dit que $A$ est un
\emph{anneau de Cohen-Macaulay} si tous les $A_{\mathfrak{p}}$ le sont. Il est
clair que si $M$ (resp. $A$) est un $A$-module de Cohen-Macaulay (resp. un
anneau de Cohen-Macaulay), $S^{-1}M$ est un $S^{-1}A$-module de Cohen-Macaulay
(resp. $S^{-1}A$ est un anneau de Cohen-Macaulay) pour toute partie
multiplicative $S$ de $A$.
\end{env}
